% Part of the grouped supplement on Boolean Doob chains.
\section{Several exceptional stationary levels}
\label{sec:several-stationary-levels}
Let $E_1,\ldots,E_N$ be disjoint nonempty subsets of the cube,
$m_i=\mu(E_i)$ and $m=\max_i m_i$. Assume $\pi=a$ off their union,
$b\le a\le1$, and $a\le\pi\le M_i$ on $E_i$. Fix $N,b,H,\kappa$
independently of dimension, with $b>0$, and assume
\[
  M_i m_i^{1/3}\le H,\qquad m_i\ge m^\kappa,\qquad
  1\le\kappa<21-8\sqrt5 . \tag{F1}
\]
The restriction compares logarithms of uniform masses. The mass and height ratios may diverge. The sets have arbitrary geometry.

\paragraph{Theorem.}
Under (F1) the spectral gap is bounded below by a positive constant
depending only on $N,b,H,\kappa$. If $\pi$ is constant on each $E_i$,
then $\Delta\le C_{N,b,\kappa}\|\pi-1\|_3$ without coordinate matching.
For arbitrary allowed heights, cubic convergence gives full recovery,
and uniform integrability of the cubed densities suffices under
diverging matching order. With a fixed $p>3$ and $\|\pi\|_p\le B$, matching
through order $k$ instead gives $\Delta\le C e^{-ck}$ with constants
depending only on $N,b,H,\kappa,p,B$; for exact levels one may take $H=B$.

For example, take disjoint sets of masses $2^{-n}$ and $2^{-2n}$ in
dimension $2n$ for $n\ge2$. Add $m_i^{-1/3}$ to a common background
on each set, choosing $a=1-\sum_i m_i^{2/3}$ for normalization.
For large $n$ the floor exceeds $1/2$ and the cubic moment stays bounded.
Both sets together violate any fixed single-spike envelope (S1):
the maximum times their combined mass to the power $1/3$ grows at
least as $2^{n/3}$. Here (F1) applies with $\kappa=2$.

\paragraph{Control between different levels.}
Set $r=12\sqrt5-15>9$, $s=(r+15)/16$, $T=P^2$, and
$\eta=1/s-1/3-\kappa/r$. Then $\eta>0$ because
$r(16/(r+15)-1/3)=21-8\sqrt5$. Differentiation shows this choice
maximizes the right side over $r>6$; this is a limit of the estimate,
not a necessary restriction for a gap. For every $\ell\ge1$, $j\ge0$,
$P^j T^\ell$ contracts $L^s$ into $L^r$.
Let $Z=\max_i\|Y\|_{r,\mu_i}$, $Y=1/g$, $\mu_i=\mu(\cdot\mid E_i)$.
The contribution from an individual level to any other satisfies
\[
 M_i\|P^jT^\ell(Y{\bf1}_{E_i})\|_{r,\mu_v}
 \le M_i m_i^{1/s}m_v^{-1/r} Z
 \le H m^\eta Z. \tag{F2}
\]
By Jensen $Y\le Pg/a$, hence
$Y\le TY+a^{-1}\sum_i M_iT(Y{\bf1}_{E_i})$.
Iterate $n=1+\lceil\log_{16}(1+\log(1/\min_i m_i))\rceil$ times
as in (S2). The unrestricted term $P^jT^nY$ has $L^{1+8\cdot16^n}$
norm at most $1/b$, so its conditional $L^r$ norm on each level is
at most $e/b$. By (F2), $Z\le e/b+(nNH/b)m^\eta Z$. For all
sufficiently small $m$, $n=O_\kappa(1+\log(1+\log(1/m)))$ and
the coefficient is at most $1/2$. Applying the iterated bound also
with $j=2$ controls $TY$ on every level. Since
$Pg=T(\pi Y)\le aTY+\sum_i M_iT(Y{\bf1}_{E_i})$, we have
\[
 \max_i\big(\|Y\|_{r,\mu_i}+\|Pg\|_{r,\mu_i}\big)\le C. \tag{F3}
\]
The conditional estimate does not require pointwise control on $g$.
When $m$ is bounded below, all the $M_i$ are bounded by (F1) and we
use the ambient $\|Y\|_9\le1/b$ instead.

\paragraph{Cut escape, including transitions between levels.}
We prove the uniform gap by contradiction along a sequence. If $m$
is bounded below, (G3) and the global $L^{9/2}$ bound on $U=\pi Y^2$
give the gap as in the one-spike proof. Pass therefore to $m\to0$.
Extract a subsequence along which every ratio $m_i/m_j$ has a limit
in $[0,\infty]$. Group indices with finite nonzero limiting ratio.
Each group $C$ denotes the union of its levels, with mass $m_C$ and
height bound $M_C=\max_{i\in C}M_i$ (indices in this maximum refer
to the group). Along this subsequence $M_Cm_C^{1/3}\le H_C<\infty$,
and (F3) holds also under $\mu_C=\mu(\cdot\mid C)$.

For $\varnothing\ne F\subset C$ write $v_F=\nu(F)$,
$\theta_F=v_F/(M_Cm_C)$ and $R_F=J(F,F)/v_F$. As in (S3) and its
cut estimate, using $\|Y\|_{8,\mu_C}\le C$ gives
\[
 R_F\le C M_C m_C^{1/3}\theta_F^{1/12}.
\]
Thus if $R_F$ stays bounded below then $\theta_F$ and
$M_Cm_C^{1/3}$ both stay bounded below, and
$v_F$ is comparable to $m_C^{2/3}$. Moreover $\nu(\cdot\mid F)$
has density at most $1/\theta_F$ relative to $\mu_C$.
Under this conditional stationary law, the expectations of $Pg$ and
$P(h{\bf1}_F)$ are uniformly bounded: for the latter,
\[
 \|P(h{\bf1}_F)\|_{3,\mu_C}
 \le m_C^{-1/3} M_C\|Y{\bf1}_C\|_{3/2}
 \le C M_C m_C^{1/3}.
\]
Also $E_{\mu_C}P{\bf1}_C\le m_C^{1/3}$ by ordinary hypercontractivity.
Consequently, for a fixed sufficiently large $A$ and all large
dimensions, at least half the conditional stationary mass has
$Pg\le A$, $P(h{\bf1}_F)\le A$ and $P{\bf1}_{F^c}\ge1/2$, by Markov's
inequality. On those vertices if $g>2A$, then $K_x(F)\le1/2$.
Otherwise the row escape bound from (S3) and its proof gives
$K_x(F^c)\ge b/(8A^2)$. Hence $R_F$ stays bounded away from one.

If the union admitted cuts $F$ with $R_F\to1$, split them into
$F_C=F\cap C$. Pass to limits of their diagonal return fractions
$R_{F_C}$, taking zero for empty pieces. The positive semidefinite
edge form gives
\[
 \sqrt{R_F}\le \sum_C \sqrt{R_{F_C}\,\nu(F_C)/\nu(F)}.
\]
Groups with vanishing diagonal return contribute $o(1)$. Of those with
positive limiting return, the one with the largest mass dominates
all the others: their cut masses are comparable to their respective
$m_C^{2/3}$, and distinct groups have mass ratios tending to zero or
infinity. The dominant summand stays bounded away from one by the
preceding row argument. This contradicts $R_F\to1$ (if there are no
such groups the sum tends to zero). The killed norm on the union is
therefore bounded away from one along the extracted subsequence by (LR1).
Off the union, $\|U\|_{9/2}$ restricted there is at most $1/b^2$.
Adding Gram multipliers as in the one-spike proof controls all small
cuts. Larger cuts obey (G3) since $E\pi^3\le1+NH^3$. Cheeger's
inequality now precludes a sequence with gap tending to zero.

\paragraph{Full spectral estimates.}
For exact levels take $M_i=\pi|_{E_i}$. Put $\delta=\|\pi-1\|_3$.
For the Lipschitz estimate it suffices to take $\delta\le1$, when
$H=2$ works. Away from levels above two, the preceding section proved
$\|\pi(Y^2-1){\bf1}_{\{\pi\le2\}}\|_3\le C_b\delta$.
On each higher level, $M_i m_i^{1/3}\le2\delta$; (F3) gives
$\|U{\bf1}_{E_i}\|_3\le C M_i m_i^{1/3}$. If $m$ is bounded below
the same conditional $L^6$ control follows from $\|Y\|_9\le1/b$
and (F1). Sum the finitely many contributions and use the Gram
inequality $\Delta\le\|U-1\|_3$. The parity example in the previous
section retains sharpness of the linear order for the full class.

\paragraph{Variable heights within levels.}
The gap statement allowed variable heights from the outset. The same
hypotheses yield a quantitative transfer for the entire spectrum even
when constant heights are unavailable. Put $\alpha=1-6/r>0$.
For any set $D$ in the small-$m$ regime, conditional H\"older on $E_i$
with $1/u=1/3-2/r$ and (F3) give
\[
 \|U{\bf1}_{E_i\cap D}\|_3
 \le C m_i^{1/3}\|\pi{\bf1}_D\|_{u,\mu_i}
 \le C (M_i m_i^{1/3})^{1-\alpha}
            \|\pi{\bf1}_{E_i\cap D}\|_3^\alpha .
\]
Here $3/u=\alpha$. Off the exceptional union use the ambient
$L^{9/2}$ bound. When $m$ stays above the absorption threshold, the
same ambient bound holds everywhere. For either regime,
\[
 \|U{\bf1}_D\|_3
 \le C_{N,b,H,\kappa}
       \big(\|\pi{\bf1}_D\|_3^\alpha+\mu(D)^{1/9}\big). \tag{F4}
\]
In particular full recovery without matching holds if
$\delta=\|\pi-1\|_3\to0$. For a rate on $\{\pi>2\}$, in the small-$m$
regime use the preceding levelwise inequality there (the background is
at most one) and $\pi\le2|\pi-1|$. On $\{\pi\le2\}$ use the inverse
estimate from the one-spike section. In the remaining, bounded-height
regime use that inverse estimate globally. This gives
$\Delta\le C(\delta+\delta^\alpha)$.

There is also a matched result at the cubic order: uniform
integrability of $\pi^3$ under $\mu$ and $k\to\infty$ imply full
spectral recovery for variable heights satisfying (F1). Indeed (F4)
transfers that uniform integrability to $U^3$. The matched bad set
$D=\{|g-1|>2^{-k/8}\}$ has exponentially small mass since the cubic
cap implies a bounded second moment. The $L^3$ norm of $U-\pi$ on $D$
therefore tends to zero; outside $D$ it is exponentially small.
To bound $QM_{\pi-1}Q$ at order three let $A$ project onto Walsh degrees
at most $\lfloor k/2\rfloor$. The low-low block vanishes by matching of
products. In each remaining block at least one high factor
$Q(I-A)$ has $L^2$ operator norm at most $2^{-k/4}$. Truncate
$\pi-1$ at a fixed height $T>1$, keeping its values on $\{\pi\le T\}$.
The bounded part in each such block costs at most $(T+1)2^{-k/4}$,
and the tail costs at most $\|\pi{\bf1}_{\{\pi>T\}}\|_3$ by
$Q:L^2\to L^3$ and the $L^2$ contractions $A,I-A$ before $Q$.
Take first the limit in $k$, then in $T$. Min--max proves the claim.
For $p>3$ and $\|\pi\|_p\le B$, the same proof is exponential:
in (F4) use $\|\pi{\bf1}_D\|_3\le B\mu(D)^{1/3-1/p}$, and take
$T=2^{k/8}$ in the spectral blocks, whose tail is at most
$B^{p/3}T^{1-p/3}$. No bound on the within-level oscillation is needed.

The \href{https://airr.science/papers/ARR-2026-2B2W1WQ7DF9N5VW7/versions/v1/}{published trapping family}
has a fixed floor and converges in cubic norm. This theorem implies
that it admits no representation satisfying (F1) with uniform
parameters. The theorem uses ordinary Boolean hypercontractivity;
\href{https://arxiv.org/abs/1606.07491}{Polyanskiy and Samorodnitsky}
give stronger estimates on small supports. Whether such an improvement
removes the logarithmic mass comparison is open here.
